\section{Quick Start: Configuring Your First Model}
\label{sec:python_quickstart}

This section describes how to install hsB--SM, prepare catchment inputs, set the parameters and initial conditions, and run the model. A forward simulation requires hillslope geometry, a channel travel-distance distribution, and meteorological time series. Observed discharge is optional for simulation and is needed for evaluation or calibration. The equations and implementation details follow in the later sections.

\subsection{Install the model}

Install Python 3.11 or newer and Git. Open a terminal and download the Python repository.

\begin{lstlisting}[language={},caption={Download the repository.}]
git clone https://github.com/mariaandradeal/1D_hsb-SM-python.git
cd 1D_hsb-SM-python
\end{lstlisting}

On Windows PowerShell, create and activate a Python environment, then install the package.

\begin{lstlisting}[language={},caption={Install on Windows.}]
python -m venv .venv
.\.venv\Scripts\Activate.ps1
python -m pip install -e .
\end{lstlisting}

If PowerShell blocks activation, use \path{.\.venv\Scripts\python.exe} in place of \texttt{python} in the installation and model commands. Activation is optional.

On macOS or Linux, use the following commands.

\begin{lstlisting}[language={},caption={Install on macOS or Linux.}]
python3 -m venv .venv
source .venv/bin/activate
python -m pip install -e .
\end{lstlisting}

\subsection{Prepare the input files}

From the repository root, copy \path{templates/catchment} to a new model directory. The copy command requires a new destination and does not overwrite an existing model. The CSV templates contain column headers only. Fill them with data for your catchment.

\begin{lstlisting}[language={},basicstyle=\ttfamily\footnotesize,caption={Create a model directory from the input template.}]
python -c "from shutil import copytree; copytree('templates/catchment', 'my_model')"
\end{lstlisting}

The model directory contains four input files. \path{geometry.csv} defines the hillslope cells, \path{channel.csv} defines the channel travel-distance distribution, \path{forcing.csv} contains meteorological and vegetation inputs, and \path{config.json} stores paths, the timestep, slopes, parameters, and initial conditions. Paths inside the configuration are relative to its directory.

In \path{geometry.csv}, supply cell-center distance \texttt{x\_m}, planform width \texttt{w\_m}, and cell length \texttt{dx\_m}, all in metres. Order cells from the outlet toward the divide. Distances, widths, and lengths must be positive. Each center must lie inside the cell edges obtained by accumulating cell lengths from the outlet. Supply prepared geometry tables; the Python package does not delineate a catchment from a digital elevation model.

In \path{channel.csv}, supply increasing positive channel distances \texttt{distance\_m} [m] and nonnegative bin weights \texttt{weight}. Their sum must be positive. The model normalizes the weights. For unequal distance-bin widths, use integrated probability masses.

The meteorological table has one row per forcing interval and the columns listed below. The forcing is spatially uniform across the hillslope.

{\small
\begin{longtable}{p{0.27\linewidth}p{0.62\linewidth}}
\toprule
Column & Meaning and unit \\
\midrule
\endhead
\texttt{date} & Start of the interval, as \texttt{YYYY-MM-DD} or \texttt{YYYY-MM-DD HH:MM:SS} \\
\texttt{P} & Precipitation rate [m\,s$^{-1}$] \\
\texttt{Temp\_C} & Air temperature [$^\circ$C] \\
\texttt{Pres\_kPa}, \texttt{VapP\_kPa} & Air pressure and actual vapor pressure [kPa] \\
\texttt{Wind\_ms} & Wind speed [m\,s$^{-1}$] \\
\texttt{SRad\_Wm2}, \texttt{LRad\_Wm2} & Incoming shortwave and longwave radiation [W\,m$^{-2}$] \\
\texttt{Albd} & Surface albedo, from 0 to 1 \\
\texttt{LAI} & Leaf area index [m$^2$\,m$^{-2}$] \\
\texttt{fv} & Vegetated area fraction, from 0 to 1 \\
\texttt{SFC} & Optional snow-cover fraction, from 0 to 1; zero if the column is absent \\
\bottomrule
\end{longtable}
}

Use one consistent time basis. Dates must be unique, increasing, and equally spaced. Required values must be finite. If precipitation is in mm/day, name its column \texttt{Prec\_mm\_day} instead of \texttt{P}; the reader converts it to m/s. Missing values are not filled automatically. To read an Excel forcing table, install \texttt{python -m pip install -e ".[excel]"} and change the forcing filename in the configuration.

\subsection{Set the configuration and initial conditions}

Open \path{my_model/config.json}. Replace the three required \texttt{null} values with numbers: \texttt{dt} is the forcing interval in seconds, \texttt{channel.slope} is the positive channel slope in m/m, and \texttt{parameters.Hill.iota\_rad} is the hillslope angle in radians. The timestep must equal the interval between forcing timestamps.

Set \texttt{forcing.lai\_basis} to \texttt{ground} when LAI is defined per unit ground area, or \texttt{vegetated} when it is defined per unit vegetated area. The reader divides ground-area LAI by \texttt{fv} once. Ground-area LAI must be zero wherever \texttt{fv} is zero.

The remaining values in the template are implementation defaults. Review them for the catchment before running. Set layer depths and hydraulic properties in \texttt{Soil.RZ} (root zone) and \texttt{Soil.TZ} (transmission zone), aquifer conductivity, drainable porosity, and thickness in \texttt{Aq}, and channel velocity and routing parameters in \texttt{Route}. Conductivities are in m\,s$^{-1}$, depths in m, and soil water contents in m$^3$\,m$^{-3}$. Van Genuchten \texttt{alpha} is in m$^{-1}$ and \texttt{n} is dimensionless. These defaults are not calibrated catchment estimates.

Choose \texttt{Flags.interception} explicitly. \texttt{Deep.phi\_d} controls the fraction of downward recharge entering the deep reservoir. Review each soil layer's initial water content, \texttt{theta0}. The optional \texttt{initial} object accepts scalar or node-vector values for \texttt{theta\_rz}, \texttt{theta\_tz}, \texttt{canopy\_m}, \texttt{aquifer\_head\_m}, \texttt{deep\_m}, and \texttt{infiltration\_m}. Soil water contents are in m$^3$\,m$^{-3}$; the other initial states are in m. Without overrides, soil water contents use \texttt{theta0}, canopy and aquifer storage start at zero, deep storage starts at 1 mm, and cumulative infiltration starts at zero. A nonzero initial canopy store requires interception to be enabled.

Six optional parameters may remain \texttt{null}. \texttt{Veg.d}, \texttt{Veg.z0m}, and \texttt{Veg.z0h} then use the canopy-height relations. \texttt{Infil.K\_mps} and \texttt{Infil.S\_mps05} then use the soil-based infiltration expressions. \texttt{Troch.alpha\_inv\_m} also defaults to the soil-based capillary length. Record data and parameter sources in \texttt{metadata} and retain an unchanged copy of the source data.

\subsection{Check the inputs and run}

After filling the tables and configuration, check that the model can load them. Resolve any reported input errors before running.

\begin{lstlisting}[language={},basicstyle=\ttfamily\footnotesize,caption={Load and check the configured inputs.}]
python -c "from hsbsm import load_config; load_config('my_model/config.json')"
\end{lstlisting}

Run the model from the repository root.

\begin{lstlisting}[language={},caption={Run with the prepared catchment inputs.}]
python -m hsbsm run my_model/config.json --output outputs/my_model
\end{lstlisting}

Use a new output directory for each experiment. Saving to an existing directory replaces files with the same names. The option \texttt{--no-plots} skips figure generation.

The same configuration can be used from a Python script.

\begin{lstlisting}[language=Python,caption={Run through the Python interface.}]
from hsbsm import load_config, run

cfg = load_config("my_model/config.json")
result = run(cfg)
result.save("outputs/my_model")
\end{lstlisting}

The notebook \path{examples/quickstart.ipynb} reads the prepared configuration and follows this workflow. It requires the user's input files. Install notebook support with \texttt{python -m pip install -e ".[notebook]"}, then start \texttt{jupyter lab}.

\subsection{Inspect the first run}

Read \path{outputs/my_model/diagnostics.json} for the water-balance residual and solver diagnostics. Inspect \path{timeseries.csv} and \path{hydrograph.png} for discharge, evapotranspiration, and soil water content. Confirm that the dates, units, and magnitudes agree with the supplied inputs. A closed water balance checks numerical accounting; evaluation of predictive performance requires observed data.

The run also exports node states and fluxes, the parameter values, monthly mean rates, annual water totals, and figures in PNG, PDF, and SVG formats. Section~\ref{sec:python_reference} defines the output conventions and numerical checks.

\newpage
